\chapter{GNSS Measurement Model}
\label{chap:gnss-measurements}

This chapter specifies the GNSS receiver measurement model of \lupnt{}: how a
terrestrial GNSS constellation is represented, how the signal-transmission
epoch is solved, which physical terms enter the pseudorange, Doppler and
carrier-phase observables, how a link is declared visible and trackable, and
what the analytic measurement Jacobian contains. The model is the one used
for the cislunar sidelobe-GNSS studies of \citet{iiyama2024navi} and
\citet{iiyama2026scud}, and follows the conventions of
\citet[Ch.~6]{iiyama2026thesis}; the underlying observation equations are the
standard ones of \citet{misra2011gnss} and \citet{kaplan2017}, with the
broadcast-message and group-delay conventions of \citet{icdgps200}.

The implementation lives in

\begin{itemize}[leftmargin=1.4em,itemsep=1pt]
  \item \file{cpp/lupnt/measurements/gnss_measurement.h} and
        \file{cpp/lupnt/measurements/gnss_measurement.cc} ---
        \cpp{GnssChannel}, \cpp{GnssMeasurement} (single channel) and
        \cpp{GNSSMeasurements} (all channels of one receiver);
  \item \file{cpp/lupnt/agents/gnss_constellation.cc} ---
        \cpp{GnssConstellation}, the PRN-indexed transmitter ephemeris and
        antenna/power metadata;
  \item \file{cpp/lupnt/interfaces/sp3_loader.cc},
        \file{cpp/lupnt/interfaces/rinex_nav_loader.cc} and
        \file{cpp/lupnt/interfaces/antex_loader.cc} --- precise, broadcast and
        antenna-phase-centre ingestion;
  \item \file{cpp/lupnt/attitude/gnss_attitude.cc} and
        \file{cpp/lupnt/attitude/gnss_yaw_steering.cc} --- the transmitter
        body frame in which the transmit-antenna gain is looked up;
  \item \file{cpp/lupnt/measurements/comms_utils.cc} and
        \file{cpp/lupnt/measurements/antenna.cc} --- the visibility test, the
        free-space path loss and the antenna pattern interpolation.
\end{itemize}

Two neighbouring chapters carry material referenced throughout: the electron
density model and the ray-traced dispersive delay are specified in
\cref{chap:ionosphere}, and the closed-form link budget, gain patterns and
tracking-loop noise mapping in \cref{chap:link-budget}. This chapter states
only how those quantities enter the GNSS observation equations.

% ===========================================================================
\section{State, Time and Frame Contract}
\label{sec:gnss-contract}

\subsection{Receiver state layout}

All vectors in this chapter are expressed in \code{options.frame}, a single
inertial frame shared by receiver and transmitter. The implementation refuses
\code{Frame::UNDEFINED}; cislunar filtering runs use \code{Frame::MOON\_CI} for
the receiver and convert the transmitter ephemeris into it, while Earth-orbit
runs leave both in \code{Frame::ECI}.

The receiver state vector is not assumed to have any particular length. It is
interpreted through the index map \cpp{GnssMeasurementOptions::indices},

\begin{equation}
  \label{eq:gnss-state-layout}
  \vec{x} =
  \begin{bmatrix}
    \cdots & \vec{r}_R\transpose & \cdots & \vec{v}_R\transpose & \cdots &
    b_R & d_R & \cdots & N & \cdots
  \end{bmatrix}\transpose ,
\end{equation}

with $\vec{r}_R \in \mathbb{R}^3$ in \si{\meter}, $\vec{v}_R \in \mathbb{R}^3$
in \si{\meter\per\second}, receiver clock bias $b_R$, clock drift $d_R$ and an
optional carrier integer ambiguity $N$ in cycles. A negative index, or an
index beyond the end of the state, disables that term: it is then neither read
nor given a Jacobian column. This is what lets the same measurement model
serve a six-state orbit-only filter, an eight-state orbit-plus-clock filter
and an ambiguity-augmented carrier-phase filter without modification.

\begin{lstlisting}[language=cpp, caption={Receiver state index map and its
defaults.}, label={lst:gnss-indices}]
// cpp/lupnt/measurements/gnss_measurement.h :: GnssMeasurementStateIndices
struct GnssMeasurementStateIndices {
  int position = 0;
  int velocity = 3;
  int clock_bias = 6;
  int clock_drift = 7;
  int carrier_integer = -1;   // < 0 disables the ambiguity state
};
\end{lstlisting}

\subsection{Clock-bias units}

The clock states are stored in \code{options.clock\_bias\_unit}, one of
\code{SECONDS}, \code{METERS} or \code{KILOMETERS}: filters that carry the
clock in range units keep the state numerically comparable to the position
components. Conversion to SI seconds happens before the range equations,

\begin{equation}
  \label{eq:gnss-clock-units}
  \Delta t_R = \operatorname{seconds}(b_R), \qquad
  \dot{\Delta t}_R = \operatorname{seconds}(d_R),
\end{equation}

and the corresponding scale factor is what appears in the clock columns of the
Jacobian (\cref{sec:gnss-jacobian}). The helpers
\cpp{ClockBiasToMeters} and \cpp{ClockBiasMetersDerivative} in
\file{cpp/lupnt/measurements/gnss_measurement.cc} implement
$c\,\operatorname{seconds}(\cdot)$ and its derivative respectively.

\subsection{Epochs and time scales}

Every epoch handed to \cpp{Compute} or \cpp{Precompute} is a receiver
\emph{signal-reception} epoch,
\begin{equation}
  \label{eq:gnss-epoch-def}
  t_R \equiv t_{\text{receive}},
\end{equation}
expressed in \code{options.receive\_time\_scale}. Constellation ephemerides
are tabulated and interpolated in $\TAI$ seconds; the implementation asserts
\code{options.ephemeris\_time\_scale == Time::TAI}. Cislunar simulations
typically set the receive scale to $\TDB$ and let the measurement model
convert to $\TAI$ for ephemeris lookup, and to $\TDB$ again for any frame
rotation. Coordinate-time to coordinate-time conversions go through
\cpp{ConvertCoordinateTime}, all others through \cpp{ConvertTime}; see
\cref{chap:time-conversions}.

Epochs are carried as \cpp{Epoch}, a split (integer second, fractional second)
value that also carries its own scale. This matters here more than anywhere
else in the library: a bare \code{double} seconds-from-J2000 epoch at
present-day dates has a unit-in-the-last-place of roughly
\SI{0.25}{\micro\second}, i.e.\ about \SI{75}{\meter} of range. The light-time
solution is therefore carried as an \emph{offset} rather than as an absolute
epoch (\cref{sec:gnss-lighttime}).

\begin{lupntnote}
  The cached-state fast path in \cpp{GnssChannel::GetTransmitState} compares
  epochs exactly. Because \cpp{Epoch} carries its scale in the value, an
  epoch with the same numeric seconds but a different scale no longer aliases
  onto the cache, which a bare \code{Real} comparison could not prevent.
\end{lupntnote}

\subsection{Frequencies}

A channel is identified by the triple (constellation, PRN, frequency). The
carrier wavelength used by every conversion between range, cycles and hertz is
$\lambda = c/f$, tabulated in \cref{tab:gnss-frequencies} together with the
spreading-code chipping rate $R_c$ that sets the code-tracking noise.

\begin{table}[htbp]
  \centering
  \caption{GNSS carrier frequencies and chipping rates, from
    \code{GNSS\_FREQ\_MAP} and \code{GNSS\_RC\_MAP} in
    \file{cpp/lupnt/devices/gnss_device.cc}. Wavelengths are
    $\lambda = c/f$ with $c = \SI{299792458}{\meter\per\second}$.}
  \label{tab:gnss-frequencies}
  \begin{tabular}{lS[table-format=4.3]S[table-format=1.6]S[table-format=2.4]}
    \toprule
    \code{GnssFreq} & {$f$ [\si{\mega\hertz}]} & {$\lambda$ [\si{\meter}]}
      & {$R_c$ [\si{\mega\hertz}]} \\
    \midrule
    \code{L1}  & 1575.420 & 0.190294 & 1.0230 \\
    \code{L2}  & 1227.600 & 0.244210 & 0.5115 \\
    \code{L5}  & 1176.450 & 0.254828 & 10.2300 \\
    \code{E1}  & 1575.420 & 0.190294 & 1.0230 \\
    \code{E6}  & 1278.750 & 0.234442 & 0.5115 \\
    \code{E5}  & 1191.795 & 0.251547 & 10.2300 \\
    \code{E5a} & 1176.450 & 0.254828 & 10.2300 \\
    \code{E5b} & 1207.140 & 0.248349 & 10.2300 \\
    \bottomrule
  \end{tabular}
\end{table}

\subsection{Symbols}

\cref{tab:gnss-symbols} fixes the notation used for the rest of the chapter.
Every quantity is a scalar unless it is typeset as a vector.

\begin{table}[htbp]
  \centering
  \caption{Symbols of the GNSS observation equations.}
  \label{tab:gnss-symbols}
  \begin{tabularx}{\textwidth}{llL}
    \toprule
    Symbol & Unit & Meaning \\
    \midrule
    $t_R$, $t_S$ & \si{\second} & signal-reception and signal-transmission epoch \\
    $\vec{r}_R,\vec{v}_R$ & \si{\meter}, \si{\meter\per\second} & receiver position/velocity at $t_R$ \\
    $\vec{r}_S,\vec{v}_S$ & \si{\meter}, \si{\meter\per\second} & transmitter position/velocity at $t_S$ \\
    $\rho$, $\dot\rho$ & \si{\meter}, \si{\meter\per\second} & geometric range and range rate \\
    $\uvec{u}$ & --- & unit line of sight, transmitter to receiver \\
    $\Delta t_R$, $\dot{\Delta t}_R$ & \si{\second}, --- & receiver clock bias and drift \\
    $\Delta t_S$, $\dot{\Delta t}_S$ & \si{\second}, --- & transmitter clock bias and drift \\
    $\Delta t_S^{\mathrm{rel}}$ & \si{\second} & transmitter periodic relativistic correction \\
    $\Delta t_S^{\mathrm{grp}}$ & \si{\second} & broadcast group delay $T_{GD}$ \\
    $\Delta\rho_{\mathrm{Sh}}$ & \si{\meter} & Shapiro (gravitational) delay \\
    $\Delta\rho_{\mathrm{pl}}$ & \si{\meter} & ionosphere/plasmasphere delay at $f$ \\
    $\Delta\rho_{\mathrm{tr}}$ & \si{\meter} & tropospheric delay (ground links only) \\
    $\lambda$, $f$ & \si{\meter}, \si{\hertz} & carrier wavelength and frequency \\
    $\phi_0$, $N$ & cycle & fractional phase bias, integer ambiguity \\
    $\varphi_{\mathrm{tx}},\theta_{\mathrm{tx}}$ & \si{\radian} & transmitter off-boresight and azimuth angle \\
    $\varphi_{\mathrm{rx}}$ & \si{\radian} & receiver off-boresight angle \\
    $C/N_0$ & \si{\dBHz} & received carrier-to-noise density ratio \\
    \bottomrule
  \end{tabularx}
\end{table}

% ===========================================================================
\section{Constellation Representation}
\label{sec:gnss-constellation}

\cpp{GnssConstellation} is a passive, receiver-independent object: it owns a
PRN-indexed transmitter ephemeris plus antenna and transmit-power metadata,
and nothing that depends on where the receiver is. Light time, visibility,
$C/N_0$ and channel selection all live in \cpp{GNSSMeasurements}, because they
require a receiver state and a receive epoch.

\subsection{PRN bookkeeping}
\label{sec:gnss-prn}

A constellation carries one \code{GnssConst} system tag (\code{GPS},
\code{GLONASS}, \code{GALILEO}, \code{BEIDOU}, \code{QZSS}) and the following
PRN-keyed containers:

\begin{table}[htbp]
  \centering
  \caption{PRN bookkeeping in \cpp{GnssConstellation}
    (\file{cpp/lupnt/agents/gnss_constellation.h}).}
  \label{tab:gnss-prn-tables}
  \begin{tabularx}{\textwidth}{lL}
    \toprule
    Member & Contents \\
    \midrule
    \code{prns\_} & ordered list of PRNs; iteration order of \cpp{BuildChannels} \\
    \code{fault\_prns\_} & PRNs declared unavailable by \cpp{SetFaultPrns}; skipped outright \\
    \code{t\_tai\_} & $[M]$ common ephemeris epoch grid, $\TAI$ seconds \\
    \code{rv\_eci\_} & PRN $\mapsto$ $[M\times 6]$ tabulated ECI state history \\
    \code{rv\_eci\_cheby\_} & PRN $\mapsto$ fitted \cpp{ChebyshevFitModel} of the same history \\
    \code{antennas\_} & PRN $\mapsto$ frequency $\mapsto$ \cpp{Antenna} gain pattern \\
    \code{P\_tx\_} & PRN $\mapsto$ frequency $\mapsto$ transmit power [\si{\dBW}] \\
    \code{freq\_list\_} & frequencies for which transmitter metadata exists \\
    \bottomrule
  \end{tabularx}
\end{table}

All PRNs of a constellation share one epoch grid \code{t\_tai\_}; a PRN whose
history has a different length is rejected at
\cpp{SetSatelliteStates}. Multiple constellations, each with its own
frequency, are attached to one receiver with
\cpp{GNSSMeasurements::AddConstellation}, and their channels are simply
concatenated, so a mixed GPS L1 plus Galileo E1 scenario is a two-entry list.
A channel therefore records \code{gnss\_const} alongside \code{prn}: PRN 1 of
GPS and PRN 1 of Galileo are distinct keys, including in the receiver's
lock-state set \code{tracking\_prns\_}, which is a set of
$(\code{GnssConst},\ \code{prn})$ pairs.

The constellation exposes three ways to become populated:
\cpp{SetSatelliteStates} (states already in memory), \cpp{LoadEphemeris}
(an HDF5 file with datasets \code{/prns}, \code{/t\_tai} and
\code{/rv\_eci/<prn>}, typically written by a Python preprocessing step), and
\cpp{SetupSatelliteStatesFromFiles} (SP3 plus ANTEX, described in
\cref{sec:gnss-sp3}). \cpp{SetupTransmitters} then attaches the gain patterns
and transmit powers.

\subsection{Piecewise Chebyshev fit of the six-dimensional state}
\label{sec:gnss-chebyshev}

Whatever the source, the runtime representation of a transmitter is a
piecewise Chebyshev model of the full six-dimensional state. For each PRN the
tabulated samples
\begin{equation}
  \label{eq:gnss-eph-table}
  \bigl\{ \bigl(t_k,\ \vec{r}_S(t_k),\ \vec{v}_S(t_k) \bigr) \bigr\}_{k=0}^{M-1},
  \qquad t_k \in \TAI ,
\end{equation}
are fitted segment by segment. On a segment centred at $t_{\mathrm{mid}}$ with
half-width $t_{\mathrm{rad}}$ the model is
\begin{equation}
  \label{eq:gnss-cheby}
  \vec{y}_S(t)
  = \begin{bmatrix} \vec{r}_S(t) \\ \vec{v}_S(t) \end{bmatrix}
  \approx \sum_{n=0}^{p} \vec{c}_n\, T_n(\tau),
  \qquad
  \tau = \frac{t - t_{\mathrm{mid}}}{t_{\mathrm{rad}}} \in [-1,1],
\end{equation}
where $T_n$ is the Chebyshev polynomial of the first kind of degree $n$ and
$\vec{c}_n \in \mathbb{R}^6$. The coefficients are obtained by sampling at the
$p+1$ Chebyshev--Gauss nodes of the segment and applying the discrete
Chebyshev transform, so \cref{eq:gnss-cheby} is the unique polynomial
interpolant of the sampler at those nodes.

Because a Chebyshev series is differentiable in closed form, the evaluator can
return the exact time derivative alongside the value,
\begin{equation}
  \label{eq:gnss-cheby-deriv}
  \frac{\dd \vec{y}_S}{\dd t}(t)
  = \frac{1}{t_{\mathrm{rad}}}\sum_{n=0}^{p} \vec{c}_n\, T_n'(\tau),
  \qquad
  T_n'(\tau) = n\,U_{n-1}(\tau),
\end{equation}
with $U_{n-1}$ the Chebyshev polynomial of the second kind. This analytic
derivative is what the SP3 loader uses to produce velocity from a
position-only fit (\cref{sec:gnss-sp3}). The constellation itself does
\emph{not} use it: it fits all six components as independent dimensions and
evaluates with the derivative output disabled, so the interpolated velocity is
the fitted tabulated velocity rather than the derivative of the fitted
position. The two agree to the accuracy of the fit but are not identical
by construction.

\begin{lstlisting}[language=cpp, caption={The constellation's state fit. Note
\code{num\_coeffs = 2} and a segment length of one tabulation step.},
label={lst:gnss-cheby-fit}]
// cpp/lupnt/agents/gnss_constellation.cc :: FitStateHistoryChebyshev
ChebyshevFitModel FitStateHistoryChebyshev(const VecXd& t_tai, const MatXd& rv_eci) {
  double t_start = t_tai(0);
  double t_end = t_tai(t_tai.size() - 1);
  double min_step = t_end - t_start;
  for (int i = 1; i < t_tai.size(); i++)
    min_step = std::min(min_step, t_tai(i) - t_tai(i - 1));

  return FitChebyshevModel(
      [&t_tai, &rv_eci](double t) {
        VecXd state(6);
        for (int k = 0; k < 6; k++) {
          VecXd col = rv_eci.col(k);
          state(k) = LinearInterp1d(t_tai, col, t);
        }
        return state;
      },
      t_start, t_end, 6, std::max(1.0, min_step), 2);
}
\end{lstlisting}

\begin{lupntwarning}
  With \code{num\_coeffs = 2} the fit is degree one, and the sampler it
  interpolates is itself a linear interpolation of the table. The
  constellation-level ``Chebyshev fit'' therefore reduces, at present, to a
  piecewise-linear interpolation of \code{rv\_eci\_} on the tabulation grid.
  It is exact at the nodes and its chord error grows quadratically with the
  tabulation step, so the ephemeris grid handed to \cpp{SetSatelliteStates}
  must be dense enough for the accuracy required --- a few tens of seconds for
  metre-level work with medium-Earth-orbit transmitters. The infrastructure
  supports higher degree unchanged; only the last two arguments of
  \cpp{FitChebyshevModel} above would change.
\end{lupntwarning}

Evaluation is a segment lookup followed by a Clenshaw recursion:

\begin{lstlisting}[language=cpp, caption={Ephemeris evaluation at an arbitrary
epoch.}, label={lst:gnss-cheby-eval}]
// cpp/lupnt/agents/gnss_constellation.cc :: GnssConstellation::GetSatelliteStateEci
Vec6 GnssConstellation::GetSatelliteStateEci(int prn, Real t_tai) const {
  auto it_cheby = rv_eci_cheby_.find(prn);
  LUPNT_CHECK(it_cheby != rv_eci_cheby_.end(),
              "Chebyshev ephemeris missing for PRN " + std::to_string(prn),
              "GnssConstellation");
  VecX state_x;
  LUPNT_CHECK(it_cheby->second.Eval(t_tai, &state_x, nullptr),
              "Ephemeris interpolation time is out of range", "GnssConstellation");
  Vec6 state = state_x;
  return state;
}
\end{lstlisting}

Queries outside $[t_0, t_{M-1}]$ throw rather than extrapolate. Since the
light-time solution moves the query about \SI{0.08}{\second} (Earth orbit) to
\SI{1.3}{\second} (lunar distance) before the first tabulated epoch, the
ephemeris window must be padded relative to the simulation window.

\subsection{Precise ephemeris ingestion (SP3)}
\label{sec:gnss-sp3}

\cpp{Sp3Loader} (\file{cpp/lupnt/interfaces/sp3_loader.cc}, exposed to Python
as \py{pnt.Sp3Loader}) parses IGS SP3 precise-ephemeris files. Satellites are
keyed by their SP3 identifiers --- \code{"G01"}, \code{"E11"}, \code{"C06"},
\code{"J03"} --- which is where the mapping from SP3 identifier to
$(\code{GnssConst},\ \code{prn})$ is made. SP3 tabulates satellite
\emph{centre-of-mass} ECEF positions and a clock bias at a fixed cadence,
typically \SI{5}{\minute} for the MGEX final products.

Interpolation is a per-satellite piecewise Chebyshev fit of position only:

\begin{itemize}[leftmargin=1.4em,itemsep=1pt]
  \item the segment length is $8$ tabulation steps, the polynomial degree is
        $\min(12, M) - 1$;
  \item the sampler feeding the fit is a local Lagrange interpolant of order
        $\min(10, M)$ --- the standard SP3 interpolation. Sampling with a
        linear interpolant instead would inject kilometre-level chord error
        over a \SI{5}{\minute} arc, which the Chebyshev fit would then
        reproduce faithfully;
  \item velocity is the analytic derivative \cref{eq:gnss-cheby-deriv} of the
        fitted position polynomial, so position and velocity are consistent by
        construction;
  \item the clock is \emph{not} fitted. It is linearly interpolated over the
        raw samples with the SP3 sentinel value
        (\num{999999.999999} \si{\micro\second}) excluded, because a Chebyshev
        fit rings at day boundaries where the sentinel appears.
\end{itemize}

\begin{lstlisting}[language=cpp, caption={Precise-ephemeris query. The returned
clock bias already contains the periodic relativistic term.},
label={lst:gnss-sp3}]
// cpp/lupnt/interfaces/sp3_loader.h :: Sp3Loader::GetPosVelClock
void Sp3Loader::GetPosVelClock(const std::string& sat_id, Real t_tai,
                               Vec6& rv_ecef, Real& clock_bias_s) const;
// cpp/lupnt/interfaces/sp3_loader.cc :: velocity from the fitted position
VecX pos, pos_dot;
it_model->second.Eval(t_tai, &pos, &pos_dot);
rv_ecef.head(3) = pos.head(3);
rv_ecef.tail(3) = pos_dot.head(3);
\end{lstlisting}

Because SP3 positions refer to the centre of mass while the signal leaves the
antenna, the antenna phase-centre offset (PCO) must be added before the state
is usable as a transmit point:
\begin{equation}
  \label{eq:gnss-pco}
  \vec{r}^{\,\mathrm{ecef}}_{\mathrm{ant}}
  = \vec{r}^{\,\mathrm{ecef}}_{\mathrm{CoM}}
  + \mat{C}_{ijk}\bigl(t, \vec{r}^{\,\mathrm{ecef}}_{\mathrm{CoM}}\bigr)\,
    \vec{\Delta}^{\,\mathrm{NEU}}_{\mathrm{PCO}} .
\end{equation}
\cpp{AntexLoader} (\file{cpp/lupnt/interfaces/antex_loader.cc}, Python
\py{pnt.AntexLoader}) parses IGS ANTEX (\code{.atx}) files, supplies
$\vec{\Delta}^{\,\mathrm{NEU}}_{\mathrm{PCO}}$ per satellite and frequency, and
applies \cref{eq:gnss-pco} through
\cpp{ApplyPcoCorrectionEcef}. The offset is a few decimetres to about two
metres and is frequency dependent, so it is not negligible at the accuracy
this model targets. The SP3 velocity is left uncorrected: the PCO is constant
in the body frame, and its contribution to ECEF velocity is below the noise of
any link-budget or visibility decision.

\begin{lupntnote}
  The PCO ``IJK'' triad, \cpp{AntexLoader::ComputeIjkToEcefRotation}
  (\code{jvec = normalize(r\_sun - r\_sat)}, \code{kvec = -normalize(r\_sat)},
  \code{ivec = jvec x kvec}), is deliberately the \emph{raw, un-orthogonalised}
  Sun-pointing frame used to generate the precomputed PCO-corrected
  ephemerides. It is a different object from the orthonormal
  \cpp{GnssAttitude} body frame used for gain-pattern lookups
  (\cref{sec:gnss-attitude}); the two must not be conflated.
\end{lupntnote}

\subsection{Broadcast ephemeris ingestion (RINEX navigation)}
\label{sec:gnss-brdc}

\cpp{RinexNavLoader} (\file{cpp/lupnt/interfaces/rinex_nav_loader.cc}, Python
\py{pnt.RinexNavLoader}) parses RINEX~V3 navigation (``BRDC'') files and
evaluates the transmitter state by Keplerian propagation of the broadcast
navigation message nearest the requested epoch, following the user algorithm
of \citet{icdgps200} and its Galileo, BeiDou and QZSS analogues. GLONASS,
whose message carries tabulated state vectors and accelerations rather than
Keplerian elements, is intentionally unsupported.

Broadcast ephemerides matter here not as a convenience but as an error source:
they carry the real-time signal-in-space error (SISE) in both position and
clock relative to the precise product, and differencing the two is how the
measurement truth acquires a realistic ephemeris/clock error.
\citet[Ch.~6.4.2 and 6.5.2]{iiyama2026thesis} treats the
phase-centre-corrected precise product as truth, derives per-satellite
broadcast errors by differencing, and removes a per-constellation median
clock-bias offset to absorb the broadcast/precise time-reference difference.

\begin{lupntnote}
  \cpp{RinexNavLoader} omits the Galileo-specific GST/GPST (``GAGP'')
  system-time-correction term. This affects the satellite \emph{clock} only,
  by well under \SI{100}{\nano\second} (below \SI{30}{\meter} of
  range-equivalent), and has no effect on the broadcast position or velocity.
\end{lupntnote}

\subsection{Synthetic constellations for future epochs}
\label{sec:gnss-almanac}

Both paths above require IGS products covering the run epoch, which do not
exist for a future mission date. For those scenarios the Lunar GNSS ODTS
driver (\code{constellation.source == "almanac"},
\file{cpp/lupnt/simulations/lunar_gnss_odts/lunar_gnss_odts_simulation.cc})
synthesises each PRN's ephemeris instead of loading it, in three steps.

\begin{enumerate}[leftmargin=1.6em,itemsep=2pt]
  \item \textbf{Seed elements.} Each PRN's coarse Keplerian set is read from a
    YUMA GPS almanac (\cpp{RinexNavLoader::LoadYumaFile}) or, absent a YUMA
    file, from the broadcast navigation message. The seed Cartesian state is
    evaluated at the almanac reference epoch ($t_k = 0$) in ECEF and rotated
    to ECI.

  \item \textbf{GPS week-number rollover.} A YUMA almanac carries a mod-1024
    week number, so its raw reference epoch can land two decades from the
    intended date. The seed epoch is snapped to the 1024-week era nearest the
    run epoch,
    \begin{equation}
      \label{eq:gnss-rollover}
      t_{\mathrm{seed}} \leftarrow t_{\mathrm{seed}}
      + \Delta_{1024}\,
        \operatorname{round}\!\left(
          \frac{t_{\mathrm{run}} - t_{\mathrm{seed}}}{\Delta_{1024}}
        \right),
      \qquad
      \Delta_{1024} = 1024 \times 7 \times \SI{86400}{\second},
    \end{equation}
    so the propagation span is the intended months rather than decades. A
    full-week BRDC seed makes this a no-op.

  \item \textbf{Numerical propagation.} From the seed state each PRN is
    propagated over the ephemeris grid under a $4\times4$ Earth gravity field
    plus Sun and Moon third-body point masses
    (\cpp{CreateGnssEarthDynamics}), and the resulting ECI states are handed to
    \cpp{GnssConstellation::SetSatelliteStates} exactly like a precise or
    broadcast fit.
\end{enumerate}

With no precise reference to difference against, the broadcast SISE must be
\emph{modelled}. The \cpp{SyntheticSISE} transmitter-error model draws a
per-PRN orbit error in the local orbital (radial, along-track, cross-track)
frame and a clock error,
\begin{equation}
  \label{eq:gnss-sise}
  \delta\vec{r}_{\mathrm{ECI}}
    = \delta_R\,\uvec{r} + \delta_A\,\uvec{t} + \delta_C\,\uvec{c},
  \qquad
  \delta_{R,A,C} \sim \mathcal{N}\!\bigl(0,\ \sigma_{R,A,C}^2\bigr),
  \qquad
  \delta t \sim \mathcal{N}\!\bigl(0,\ \sigma_{\mathrm{clk}}^2\bigr),
\end{equation}
once per PRN from a seeded random number generator, held fixed over the (short)
arc --- a first-order stand-in for the near-constant systematic error a real
receiver sees over a pass. The magnitudes come from the four configuration
keys \code{constellation.synthetic\_sise\_radial\_m},
\code{\ldots\_along\_m}, \code{\ldots\_cross\_m} and \code{\ldots\_clock\_m}
(defaults \SI{0.3}{\meter}, \SI{0.8}{\meter}, \SI{0.8}{\meter} and
\SI{0.5}{\meter}), and the delta is injected through the same
filter-transmitter code path (\cpp{BroadcastEphemerisError}) that the measured
broadcast error uses.

% ===========================================================================
\section{Light Time and Signal-Transmit-Time Iteration}
\label{sec:gnss-lighttime}

The observables relate a receiver state at $t_R$ to a transmitter state at the
earlier epoch $t_S$ at which the signal left the antenna. In an inertial frame
the transmit epoch is the root of
\begin{equation}
  \label{eq:gnss-lighttime-root}
  t_R - t_S - \frac{\bigl\lVert \vec{r}_R(t_R) - \vec{r}_S(t_S) \bigr\rVert}{c} = 0 ,
\end{equation}
solved by the fixed-point iteration
\begin{equation}
  \label{eq:gnss-lighttime-fixed}
  t_S^{(0)} = t_R, \qquad
  t_S^{(i+1)} = t_R
    - \frac{\bigl\lVert \vec{r}_R(t_R) - \vec{r}_S\bigl(t_S^{(i)}\bigr) \bigr\rVert}{c} .
\end{equation}
The map is a contraction with factor $\lVert \vec{v}_S \rVert / c \approx
\num{1.3e-5}$ for a medium-Earth-orbit transmitter, so two iterations already
reach the \SI{1e-11}{\second} default tolerance. The iteration stops when
\begin{equation}
  \label{eq:gnss-lighttime-stop}
  \bigl| t_S^{(i+1)} - t_S^{(i)} \bigr| < \epsilon_{\mathrm{lt}},
  \qquad \epsilon_{\mathrm{lt}} = \code{options.light\_time\_tolerance\_s},
\end{equation}
or after at most
\code{options.light\_time\_max\_iterations} passes
(default $10$). Setting \code{options.solve\_light\_time} to false freezes
$t_S = t_R$, which is useful only for comparing against models that ignore
light time.

\subsection{The offset formulation}

\cref{eq:gnss-lighttime-fixed} is never evaluated as written. Subtracting
$\rho/c$ from an absolute epoch of order \num{8e8} seconds discards roughly
\SI{1e-7}{\second} --- about \SI{30}{\meter} of light time --- to the double
unit-in-the-last-place. The implementation instead iterates on the
\emph{offset}
\begin{equation}
  \label{eq:gnss-lighttime-offset}
  \Delta\tau^{(i+1)} = -\frac{\bigl\lVert \vec{r}_R - \vec{r}_S\bigl(t_R + \Delta\tau^{(i)}\bigr) \bigr\rVert}{c},
  \qquad
  t_S = t_R + \Delta\tau ,
\end{equation}
holding the absolute receive epoch fixed. Carried on an \cpp{Epoch}, the
offset lives entirely in the fractional-second field and is never absorbed
into the epoch's unit-in-the-last-place. The absolute transmit epoch is formed
only to \emph{evaluate} the ephemeris, where a residual
\SI{0.25}{\micro\second} maps to sub-millimetre position error at
\si{\kilo\meter\per\second} transmitter speeds.

\begin{lstlisting}[language=cpp, caption={Light-time iteration on the offset,
inside the per-PRN channel builder.}, label={lst:gnss-lighttime}]
// cpp/lupnt/measurements/gnss_measurement.cc :: GNSSMeasurements::BuildChannelForPrn
Epoch t_rx_ephem = Epoch::FromSeconds(receive_time, options_.receive_time_scale)
                       .To(options_.ephemeris_time_scale);
Real dtau = Real(0.0);
Epoch t_tx_ephem = t_rx_ephem;
Vec6 rv_tx = ConvertEciTransmitStateToMeasurementFrame(
    constellation->GetSatelliteStateEci(prn, t_tx_ephem.ToSeconds()),
    t_tx_ephem.ToSeconds(), options_);

if (options_.solve_light_time) {
  for (int iter = 0; iter < options_.light_time_max_iterations; iter++) {
    Real rho = (r_rx - rv_tx.head(3)).norm();
    Real dtau_next = -rho / Real(C);
    bool converged = abs(dtau_next - dtau) < options_.light_time_tolerance_s;
    dtau = dtau_next;
    t_tx_ephem = t_rx_ephem + dtau;
    rv_tx = ConvertEciTransmitStateToMeasurementFrame(
        constellation->GetSatelliteStateEci(prn, t_tx_ephem.ToSeconds()),
        t_tx_ephem.ToSeconds(), options_);
    if (converged) break;
  }
}
\end{lstlisting}

Clock offsets, group delay and the propagation delays are \emph{not} used to
solve \cref{eq:gnss-lighttime-root}: the iteration solves the physical
light-time geometry, and every clock and medium term enters additively in the
observation equations of \cref{sec:gnss-observables}. This separation is what
makes the model linear in the clock states.

\subsection{Truth generation versus filter evaluation}

A precomputed channel stores its converged \code{transmit\_time} and
\code{tx\_state}. When the filter later evaluates the same channel at a
\emph{different} receiver state, the stored transmit epoch is no longer the
light-time solution for that state. The flag
\code{options.recompute\_transmit\_time\_from\_ephemeris} selects the
behaviour:

\begin{itemize}[leftmargin=1.4em,itemsep=1pt]
  \item \code{false} (default) --- reuse the stored transmit epoch and state.
    Correct for truth generation, where the light-time solution belongs to the
    truth trajectory.
  \item \code{true} --- re-solve \cref{eq:gnss-lighttime-offset} from the
    current receiver position against the channel's own Chebyshev coefficients
    (\cpp{ResolveTransmitState}). Correct for estimated measurements, where
    the predicted observable must be a consistent function of the estimated
    state.
\end{itemize}

The difference is second order --- the transmit-epoch shift is
$\uvec{u}\transpose\delta\vec{r}_R/c$, moving the transmitter by
$\lVert\vec{v}_S\rVert/c \approx \num{1.3e-5}$ of the position error --- but it
matters when the same channel is reused across a long filter update.

% ===========================================================================
\section{Geometry}
\label{sec:gnss-geometry}

After light-time convergence, write $\vec{r}_S = \vec{r}_S(t_S)$ and
$\vec{v}_S = \vec{v}_S(t_S)$, and define
\begin{align}
  \label{eq:gnss-delta}
  \Delta\vec{r} &= \vec{r}_R - \vec{r}_S, &
  \Delta\vec{v} &= \vec{v}_R - \vec{v}_S, \\
  \label{eq:gnss-range}
  \rho &= \lVert \Delta\vec{r} \rVert, &
  \uvec{u} &= \frac{\Delta\vec{r}}{\rho}, \\
  \label{eq:gnss-rangerate}
  \dot\rho &= \frac{\Delta\vec{r}\transpose \Delta\vec{v}}{\rho}
            = \uvec{u}\transpose\Delta\vec{v} . & &
\end{align}

The sign convention is fixed once here and used everywhere below:
\emph{$\uvec{u}$ points from the transmitter to the receiver}, so $\dot\rho >
0$ means the link is opening (the receiver is receding) and the received
frequency is \emph{lower} than the transmitted one. The Doppler observable of
\cref{sec:gnss-observables} carries the compensating minus sign.

\begin{figure}[htbp]
  \centering
  \begin{tikzpicture}[x=1.05cm, y=1.05cm, >=Latex]

    % --- occluding body (Earth) -------------------------------------------
    \fill[black!8] (0,0) circle (1.15);
    \draw[thick, black!45] (0,0) circle (1.15);
    \node[font=\small\sffamily, black!55] at (0.42,-0.72) {Earth};
    \draw[black!45, dashed] (0,0) -- (-0.81,0.81);
    \node[font=\scriptsize, black!55] at (-0.28,0.62) {$R_\oplus$};

    % --- receiver body (Moon) ---------------------------------------------
    \fill[black!8] (5.9,1.95) circle (0.42);
    \draw[thick, black!45] (5.9,1.95) circle (0.42);
    \node[font=\scriptsize\sffamily, black!55] at (5.9,1.95) {Moon};

    % --- transmitter at t_S and at t_R ------------------------------------
    \coordinate (TS) at (-4.35,1.55);
    \coordinate (TR) at (-4.05,2.05);
    \coordinate (RX) at (4.75,1.15);
    \coordinate (TB) at (-4.15,-1.85);

    % visible link
    \draw[model] (TS) -- (RX);
    \node[font=\scriptsize, above, sloped] at ($ (TS)!0.42!(RX) $)
      {$\rho = c\,(t_R - t_S)$};

    % transmitter motion during the light time
    \draw[flow, black!55] (TS) -- (TR);
    \node[font=\scriptsize, black!55, above right=-1pt and 1pt] at (TR)
      {$\vec{r}_S(t_R)$};

    % nadir / boresight of the transmitter
    \draw[flow, blue!65!black] (TS) -- ($ (TS)!0.55!(0,0) $);
    \node[font=\scriptsize, blue!65!black, below left=0pt and 1pt]
      at ($ (TS)!0.55!(0,0) $) {$\uvec{e}_z$};
    \draw[blue!65!black] ($ (TS)!0.9cm!(0,0) $) arc (-9:22:0.9);
    \node[font=\scriptsize, blue!65!black] at ($ (TS)+(1.15,0.28) $)
      {$\varphi_{\mathrm{tx}}$};

    % occulted link
    \draw[integral, red!75!black, dashed] (TB) -- (RX);
    \node[font=\scriptsize, red!75!black, below right, sloped]
      at ($ (TB)!0.78!(RX) $) {occulted};

    % grazing / tangent ray from the blocked transmitter
    \draw[black!45, dotted] (TB) -- ($ (TB)!1.5!(-0.72,0.90) $);
    \node[font=\scriptsize, black!55, right=1pt]
      at ($ (TB)!1.5!(-0.72,0.90) $) {grazing ray};

    % receiver boresight
    \draw[flow, orange!85!black] (RX) -- ($ (RX)!0.75!(5.9,1.95) $);
    \node[font=\scriptsize, orange!85!black, below right=0pt and 1pt]
      at ($ (RX)!0.75!(5.9,1.95) $) {$\uvec{b}_R$};
    \node[font=\scriptsize, orange!85!black] at ($ (RX)+(-0.15,0.75) $)
      {$\varphi_{\mathrm{rx}}$};

    % --- nodes -------------------------------------------------------------
    \node[boxnode, minimum width=2.0cm, anchor=east] at ($ (TS)+(0.05,0) $)
      {GNSS SV\\[-1pt]{\scriptsize $\vec{r}_S(t_S),\vec{v}_S(t_S)$}};
    \node[boxnode, minimum width=2.0cm, anchor=east] at ($ (TB)+(0.05,0) $)
      {GNSS SV\\[-1pt]{\scriptsize blocked PRN}};
    \node[boxnode, minimum width=1.9cm, anchor=west] at ($ (RX)+(-0.05,-0.95) $)
      {receiver\\[-1pt]{\scriptsize $\vec{r}_R(t_R),\vec{v}_R(t_R)$}};

    \fill (TS) circle (1.6pt);
    \fill[black!45] (TR) circle (1.2pt);
    \fill (TB) circle (1.6pt);
    \fill (RX) circle (1.6pt);

  \end{tikzpicture}
  \caption{Transmit/receive geometry. The signal leaves the transmitter at
    $t_S$ and arrives at $t_R$; during the light time the transmitter moves to
    $\vec{r}_S(t_R)$, and it is the \emph{retarded} state that enters the
    observation equations. The transmit antenna gain is evaluated at the
    off-boresight angle $\varphi_{\mathrm{tx}}$ from the nadir axis
    $\uvec{e}_z$ of the yaw-steering body frame, and the receive gain at
    $\varphi_{\mathrm{rx}}$ from the receiver boresight $\uvec{b}_R$. The
    lower PRN is rejected because the Earth blocks the line of sight.}
  \label{fig:gnss-geometry}
\end{figure}

\begin{lupntnote}
  All three of $\vec{r}_R$, $\vec{v}_R$ and $\vec{r}_S$, $\vec{v}_S$ are in
  the same inertial frame, so \cref{eq:gnss-rangerate} is the full inertial
  range rate. No Sagnac term appears, because none is needed: the Sagnac
  correction is exactly the difference between an Earth-fixed and an inertial
  formulation, and this model never leaves the inertial frame.
\end{lupntnote}

% ===========================================================================
\section{Clock Terms}
\label{sec:gnss-clocks}

\subsection{Receiver clock}

The receiver clock bias and drift enter as a range and a range-rate offset,
\begin{equation}
  \label{eq:gnss-rx-clock}
  \Delta\rho_R = c\,\Delta t_R, \qquad
  \Delta\dot\rho_R = c\,\dot{\Delta t}_R ,
\end{equation}
with $\Delta t_R > 0$ meaning the receiver clock \emph{leads} true time, which
makes the measured pseudorange longer than the geometric range. This is the
convention of \citet{misra2011gnss}: the receiver clock term is added.

\subsection{Transmitter clock}

The transmitter contributes three separate effects, combined into one
effective bias:
\begin{equation}
  \label{eq:gnss-tx-clock}
  \Delta t_S^{\mathrm{eff}}
  = \Delta t_S + \Delta t_S^{\mathrm{rel}} - \Delta t_S^{\mathrm{grp}} ,
  \qquad
  \dot{\Delta t}_S^{\mathrm{eff}} = \dot{\Delta t}_S ,
\end{equation}
and the corresponding range terms are
\begin{equation}
  \label{eq:gnss-tx-clock-range}
  \Delta\rho_S = c\,\Delta t_S^{\mathrm{eff}}, \qquad
  \Delta\dot\rho_S = c\,\dot{\Delta t}_S^{\mathrm{eff}} .
\end{equation}
The transmitter clock term is \emph{subtracted} from the observables: a
satellite clock that leads true time transmits early, so the signal arrives
having apparently travelled a shorter distance.

\begin{description}[leftmargin=1.4em,style=nextline,font=\normalfont\itshape]
  \item[$\Delta t_S$ --- \code{tx\_clock\_bias\_s}]
    The satellite clock offset. From SP3 it is the interpolated precise clock
    (which, in \cpp{Sp3Loader}, already contains the relativistic term); from
    BRDC it is the polynomial $a_{f0} + a_{f1}(t-t_{oc}) + a_{f2}(t-t_{oc})^2$
    of \citet{icdgps200}.
  \item[$\Delta t_S^{\mathrm{rel}}$ --- \code{relativistic\_correction\_s}]
    The periodic relativistic correction due to orbital eccentricity, in the
    state-based broadcast form
    \begin{equation}
      \label{eq:gnss-relativity}
      \Delta t_S^{\mathrm{rel}}
      = -\frac{2\,\vec{r}_S\transpose\vec{v}_S}{c^{2}} ,
    \end{equation}
    which is the Cartesian equivalent of the broadcast
    $-2\sqrt{\mu}\,e\sqrt{A}\sin E / c^2$. Its amplitude is
    $2 e \sqrt{\mu A}/c^2$, about \SI{23}{\nano\second} (\SI{7}{\meter}) for a
    GPS satellite with $e = 0.01$. Enabled by
    \code{options.apply\_transmitter\_relativity} (default true).
  \item[$\Delta t_S^{\mathrm{grp}}$ --- \code{group\_delay\_s}]
    The broadcast group delay $T_{GD}$ (or the constellation's BGD/TGD
    equivalent), which calibrates the satellite clock, defined on the
    ionosphere-free combination, to the single frequency actually tracked
    \citep{icdgps200}. It is subtracted from the effective bias, hence added
    to the pseudorange. \lupnt{} carries the field but does not populate it
    from the navigation message; it defaults to zero.
\end{description}

\begin{lstlisting}[language=cpp, caption={Effective transmitter clock terms.},
label={lst:gnss-tx-clock}]
// cpp/lupnt/measurements/gnss_measurement.cc :: GnssChannel
Real GnssChannel::EffectiveTransmitClockBiasSeconds() const {
  return tx_clock_bias_s + relativistic_correction_s - group_delay_s;
}
Real GnssChannel::EffectiveTransmitClockDrift() const { return tx_clock_drift; }

// cpp/lupnt/measurements/gnss_measurement.cc :: anonymous namespace
Real TransmitterRelativisticCorrectionSeconds(const Vec6& rv_tx) {
  return -2.0 * rv_tx.head(3).dot(rv_tx.tail(3)) / (Real(C) * Real(C));
}
\end{lstlisting}

\begin{lupntwarning}
  \cref{eq:gnss-relativity} must be evaluated in the frame in which the
  satellite clock is defined, i.e.\ Earth-centred inertial. The
  implementation applies it to \code{rv\_tx} \emph{after} conversion into
  \code{options.frame}. For \code{Frame::ECI} the two coincide; for a
  Moon-centred frame the product $\vec{r}_S\transpose\vec{v}_S$ is not the
  same quantity, and the correction is wrong by the difference. Runs that use
  a lunar measurement frame and also require the periodic relativistic term at
  the nanosecond level should supply the correction externally through
  \code{relativistic\_correction\_s} with
  \code{apply\_transmitter\_relativity = false}.
\end{lupntwarning}

% ===========================================================================
\section{Propagation Delays}
\label{sec:gnss-propagation}

\subsection{Shapiro delay}
\label{sec:gnss-shapiro}

When \code{options.apply\_shapiro\_delay} is true (the default), the one-way
gravitational (Shapiro) delay of the Sun is added as a range correction,
\begin{equation}
  \label{eq:gnss-shapiro}
  \Delta\rho_{\mathrm{Sh}}
  = \frac{2\mu_\odot}{c^{2}}
    \ln\!\left(
      \frac{r_R^{\odot} + r_S^{\odot} + \rho}
           {r_R^{\odot} + r_S^{\odot} - \rho}
    \right),
  \qquad
  \begin{aligned}
    r_R^{\odot} &= \lVert \vec{r}_R - \vec{r}_\odot(t_R) \rVert, \\
    r_S^{\odot} &= \lVert \vec{r}_S - \vec{r}_\odot(t_R) \rVert .
  \end{aligned}
\end{equation}
It is non-dispersive and always positive, so it is added to \emph{both} the
code and the carrier observable. The Sun position $\vec{r}_\odot(t_R)$ comes
from \cpp{SetSunPositionProvider}; with no provider installed the model falls
back to a fixed placeholder $(0, \SI{1}{\au}, 0)$, which is adequate for the
magnitude of the term but not for anything else. For a cislunar link
$2\mu_\odot/c^2 \approx \SI{2.95}{\kilo\meter}$ multiplies a logarithm of
order \num{1e-5}, giving a few centimetres.

\begin{lstlisting}[language=cpp, caption={Shapiro delay, guarded against the
degenerate geometry where the log argument is non-positive.},
label={lst:gnss-shapiro}]
// cpp/lupnt/measurements/gnss_measurement.cc :: GNSSMeasurements::ComputeShapiroDelay
Vec3 r_sun = SunPosition(receive_time);
Real rx_norm = (r_rx_eci - r_sun).norm();
Real tx_norm = (r_tx_eci - r_sun).norm();
Real rho = (r_rx_eci - r_tx_eci).norm();
if (rx_norm <= 0.0 || tx_norm <= 0.0) return Real(0.0);

Real numerator = rx_norm + tx_norm + rho;
Real denominator = rx_norm + tx_norm - rho;
if (denominator <= 0.0 || numerator <= 0.0) return Real(0.0);

return 2.0 * Real(GM_SUN) / (Real(C) * Real(C)) * log(numerator / denominator);
\end{lstlisting}

Only the Sun is modelled. The \code{shapiro\_mu} and
\code{shapiro\_body\_position} option fields are deprecated placeholders and
are not read. Earth's contribution to a cislunar link is a few millimetres and
is omitted.

\subsection{Ionospheric and plasmaspheric delay}
\label{sec:gnss-plasma}

The dispersive delay of the terrestrial ionosphere and plasmasphere is the
dominant medium error for a lunar receiver looking back through the Earth's
plasma envelope. It is represented as a positive \emph{code} delay in metres,
\begin{equation}
  \label{eq:gnss-plasma}
  \Delta\rho_{\mathrm{pl}}
  = \frac{40.3}{f^{2}} \int_\Gamma n_e \,\dd s \;\ge\; 0 ,
\end{equation}
with $n_e$ the electron density along the (possibly bent) ray path $\Gamma$,
$f$ in \si{\hertz} and $\Delta\rho_{\mathrm{pl}}$ in \si{\meter} when the
integral is in \si{\per\meter\squared}. Being dispersive, it advances the
carrier by as much as it delays the code:
\begin{equation}
  \label{eq:gnss-plasma-sign}
  \text{code observable uses } +\Delta\rho_{\mathrm{pl}},
  \qquad
  \text{carrier observable uses } -\Delta\rho_{\mathrm{pl}} .
\end{equation}
This is the code--carrier divergence that a dual-frequency or
carrier-smoothing algorithm exploits.

The term is \emph{disabled by default}
(\code{options.apply\_ionosphere\_plasma\_delay = false}). When enabled,
\cpp{ComputeIonospherePlasmaDelay} dispatches in strict priority order:

\begin{enumerate}[leftmargin=1.6em,itemsep=1pt]
  \item \cpp{SetCustomIonospherePlasmaDelayModel} --- a per-link callback
    $(t_R, \vec{r}_R, \vec{r}_S, f) \mapsto \Delta\rho_{\mathrm{pl}}$;
  \item \cpp{SetIonospherePlasmaRayTraceOptions} --- the built-in
    GCPM/IRI ray tracer in
    \file{cpp/lupnt/environment/plasma/tec/raytrace.cc}, returning either the
    TEC-only delay or TEC plus the second- and third-order terms according to
    \code{delay\_mode};
  \item \code{options.default\_ionosphere\_plasma\_delay\_m} --- a constant.
\end{enumerate}

The batch variant \cpp{SetBatchCustomIonospherePlasmaDelayModel} is evaluated
once for all epochs and channels; see \cref{sec:gnss-precompute}. The
electron-density model, the ray-bending integration and the higher-order terms
are specified in full in \cref{chap:ionosphere}; this chapter fixes only
\cref{eq:gnss-plasma-sign}, the interface, and the unit (metres of code delay
at the channel frequency).

\begin{lupntnote}
  The ray tracer works in geocentric Earth-fixed coordinates
  (\code{raytrace\_frame = Frame::ECEF}) and expects $\UTC$ seconds from J2000
  (\code{raytrace\_epoch\_scale = Time::UTC}), so the measurement model
  converts both endpoints and the epoch before calling it, and converts
  nothing back --- the delay is a scalar. Positions are passed in
  \si{\kilo\meter}.
\end{lupntnote}

\subsection{Tropospheric delay}
\label{sec:gnss-tropo}

The troposphere is non-dispersive and is added to both observables with the
same sign as the Shapiro term. \lupnt{} implements the Saastamoinen zenith
delay with a $1/\sin e$ mapping \citep{saastamoinen1972},
\begin{equation}
  \label{eq:gnss-tropo}
  \Delta\rho_{\mathrm{tr}}
  = \frac{1}{\sin e}
    \left[
      \frac{\num{0.0022768}\,P}{f(\varphi, h)}
      + \frac{\num{0.0022768}}{f(\varphi, h)}
        \left( \frac{1255}{T} + 0.05 \right) e_w
    \right],
\end{equation}
with $e$ the elevation angle, $P$ the total pressure in \si{\hecto\pascal},
$T$ the temperature in \si{\kelvin}, $e_w$ the water-vapour partial pressure
from a Tetens saturation formula, and
$f(\varphi,h) = 1 - \num{0.00266}\cos 2\varphi - \num{2.8e-7} h$ the
latitude/height gravity factor. The horizon singularity is guarded by clamping
$\sin e \ge \num{1e-3}$.

\begin{lupntwarning}
  \cref{eq:gnss-tropo} lives in
  \file{cpp/lupnt/measurements/ground_station_corrections.cc}
  (\cpp{TroposphereDelaySaastamoinen}) and is used by the ground-station
  tracking application, \emph{not} by \cpp{GNSSMeasurements}. The GNSS channel
  has no troposphere field: the model targets space receivers, whose links do
  not traverse the neutral atmosphere. A ground-based GNSS receiver simulated
  with this model will be missing \SIrange{2.3}{25}{\meter} of delay
  depending on elevation.
\end{lupntwarning}

\subsection{Phase wind-up}
\label{sec:gnss-windup}

A right-hand circularly polarised carrier accumulates one cycle of phase per
relative rotation of the transmit and receive antennas about the line of
sight. The standard model is
\begin{equation}
  \label{eq:gnss-windup}
  \Delta\phi_{\mathrm{wu}}
  = \operatorname{sign}\!\bigl(\uvec{u}\transpose(\uvec{D}_S\times\uvec{D}_R)\bigr)
    \cos^{-1}\!\left(
      \frac{\uvec{D}_S\transpose\uvec{D}_R}
           {\lVert \uvec{D}_S\rVert\,\lVert \uvec{D}_R\rVert}
    \right)
  \quad\text{[cycle]},
\end{equation}
with $\uvec{D}_S$ and $\uvec{D}_R$ the effective dipole vectors of the two
antennas projected onto the plane normal to $\uvec{u}$, unwrapped in time
against the previous epoch. \lupnt{} does \emph{not} implement
\cref{eq:gnss-windup}: there is no wind-up term anywhere in the measurement
code. Its effect is absorbed into \code{phase\_bias\_cycles} and, over an arc,
into the estimated ambiguity, which is acceptable while the receiver attitude
is slowly varying and no cycle-slip reset occurs. It is listed here because it
is the largest \emph{unmodelled} carrier-phase term for a spinning or
rapidly slewing receiver --- up to one full cycle, i.e.\ \SI{19}{\centi\meter}
at L1, per revolution.

% ===========================================================================
\section{Observation Equations}
\label{sec:gnss-observables}

Collecting the pieces, with $\lambda = c/f$ and the clock range terms of
\cref{eq:gnss-rx-clock,eq:gnss-tx-clock-range}, the three observables are as
follows. Each is emitted only if its \code{GnssObservable} appears in
\code{options.observables}; the default is all three, in the order
pseudorange, Doppler, carrier phase.

\subsection{Pseudorange}

\begin{equation}
  \label{eq:gnss-pseudorange}
  \boxed{\;
  P = \rho
    + \Delta\rho_{\mathrm{Sh}}
    + \Delta\rho_{\mathrm{pl}}
    + \Delta\rho_R
    - \Delta\rho_S \;}
  \qquad [\si{\meter}]
\end{equation}

Every term is in metres. The geometric range $\rho$ is the light-time-corrected
range \cref{eq:gnss-range}; the Shapiro and plasma delays are added because
both retard the code; the receiver clock is added and the effective
transmitter clock (bias plus relativity minus group delay) subtracted per
\cref{sec:gnss-clocks}.

\subsection{Doppler}

\begin{equation}
  \label{eq:gnss-doppler}
  \boxed{\;
  D = -\frac{\dot\rho + \Delta\dot\rho_R - \Delta\dot\rho_S}{\lambda} \;}
  \qquad [\si{\hertz}]
\end{equation}

The leading minus sign follows from the $\uvec{u}$ convention of
\cref{sec:gnss-geometry}: an opening link ($\dot\rho > 0$) produces a negative
frequency shift. $D$ is the shift of the received carrier relative to the
nominal frequency $f$, not the received frequency itself. Multiplying by
$-\lambda$ recovers the pseudorange rate in \si{\meter\per\second}, which is
the form the Endurance-style single-satellite Doppler navigation problem uses.

\subsection{Carrier phase}

\begin{equation}
  \label{eq:gnss-carrier}
  \boxed{\;
  \Phi = \frac{\rho
    + \Delta\rho_{\mathrm{Sh}}
    - \Delta\rho_{\mathrm{pl}}
    + \Delta\rho_R
    - \Delta\rho_S}{\lambda}
    + \phi_0 \;}
  \qquad [\text{cycle}]
\end{equation}

with $\phi_0 = \code{phase\_bias\_cycles}$ the fractional (non-integer) bias
lumping the transmit and receive hardware phase offsets and, in this model,
the unmodelled wind-up. Note the sign flip on $\Delta\rho_{\mathrm{pl}}$
relative to \cref{eq:gnss-pseudorange}. When an integer ambiguity is present
--- either in the channel (\code{integer\_ambiguity\_cycles}) or, taking
precedence, in the state at \code{indices.carrier\_integer} --- the emitted
observable is
\begin{equation}
  \label{eq:gnss-carrier-meas}
  \Phi_{\mathrm{meas}} = \Phi + N .
\end{equation}

\cref{tab:gnss-signs} summarises the sign convention for every term, which is
the single most error-prone part of any GNSS implementation.

\begin{table}[htbp]
  \centering
  \caption{Sign of each term in the three observables. ``$+$'' means the term
    increases the observable as written; entries in the Doppler column refer
    to the bracketed range-rate sum of \cref{eq:gnss-doppler}, before the
    overall $-1/\lambda$.}
  \label{tab:gnss-signs}
  \begin{tabular}{llccc}
    \toprule
    Term & Field & $P$ & $[\,\cdot\,]$ of $D$ & $\Phi$ \\
    \midrule
    geometric range / rate   & ---                             & $+\rho$ & $+\dot\rho$ & $+\rho/\lambda$ \\
    receiver clock           & state \code{clock\_bias/drift}  & $+$ & $+$ & $+$ \\
    transmitter clock        & \code{tx\_clock\_bias\_s}       & $-$ & $-$ & $-$ \\
    transmitter relativity   & \code{relativistic\_correction\_s} & $-$ & $\cdot$ & $-$ \\
    group delay              & \code{group\_delay\_s}          & $+$ & $\cdot$ & $+$ \\
    Shapiro                  & \code{shapiro\_delay\_m}        & $+$ & $\cdot$ & $+$ \\
    ionosphere / plasma      & \code{ionosphere\_plasma\_delay\_m} & $+$ & $\cdot$ & $-$ \\
    fractional phase bias    & \code{phase\_bias\_cycles}      & $\cdot$ & $\cdot$ & $+$ \\
    integer ambiguity        & \code{integer\_ambiguity\_cycles} & $\cdot$ & $\cdot$ & $+$ \\
    \bottomrule
  \end{tabular}
\end{table}

The whole of \cref{eq:gnss-pseudorange,eq:gnss-doppler,eq:gnss-carrier} is
sixteen lines of code:

\begin{lstlisting}[language=cpp, caption={The three observables.},
label={lst:gnss-observables}]
// cpp/lupnt/measurements/gnss_measurement.cc :: GnssMeasurement::ComputeValue
Vec3 dr = r_rx - r_tx;
Vec3 dv = v_rx - v_tx;
Real rho = dr.norm();
Real rho_dot = dr.dot(dv) / rho;

Real tx_bias_m  = C * channel_.EffectiveTransmitClockBiasSeconds();
Real tx_drift_mps = C * channel_.EffectiveTransmitClockDrift();
Real code_delay_m    = channel_.shapiro_delay_m + channel_.ionosphere_plasma_delay_m;
Real carrier_delay_m = channel_.shapiro_delay_m - channel_.ionosphere_plasma_delay_m;
Real range_with_clock      = rho + code_delay_m + rx_bias_m - tx_bias_m;
Real range_rate_with_clock = rho_dot + rx_drift_mps - tx_drift_mps;
Real lambda = channel_.Wavelength();

value.pseudorange_m = range_with_clock;
value.doppler_hz = -range_rate_with_clock / lambda;
value.carrier_phase_cycles
    = (rho + carrier_delay_m + rx_bias_m - tx_bias_m) / lambda + channel_.phase_bias_cycles;
value.carrier_integer_cycles = integer_cycles;
\end{lstlisting}

\begin{lupntnote}
  The Doppler observable contains no time derivative of the Shapiro or plasma
  delay. $\dd\Delta\rho_{\mathrm{Sh}}/\dd t$ is of order
  \SI{1e-5}{\meter\per\second} and negligible;
  $\dd\Delta\rho_{\mathrm{pl}}/\dd t$ can reach
  \SIrange{0.01}{0.1}{\meter\per\second} at low geocentric altitude and is a
  genuine omission for Doppler-only navigation through the plasmasphere.
\end{lupntnote}

\subsection{Time-differenced carrier phase}
\label{sec:gnss-tdcp}

The Lunar GNSS ODTS scenario also forms a time-differenced carrier phase
(TDCP) observable, expressed as a carrier \emph{range} difference in metres
between consecutive epochs,
\begin{equation}
  \label{eq:gnss-tdcp}
  \Delta\Phi_{k,k-1}
  = \lambda\bigl(\Phi_k - \Phi_{k-1}\bigr)
  = \bigl(\rho_k - \rho_{k-1}\bigr)
  + \bigl(\Delta\rho_{R,k} - \Delta\rho_{R,k-1}\bigr)
  - \bigl(\Delta\rho_{S,k} - \Delta\rho_{S,k-1}\bigr)
  + \cdots ,
\end{equation}
in which the ambiguity $N$ and the fractional bias $\phi_0$ cancel exactly as
long as lock is maintained. Because \cref{eq:gnss-tdcp} depends on both the
current and the previous receiver state, it cannot be processed by a plain
Kalman update; \lupnt{} processes it with the UDU stochastic-cloning filter
\citep{iiyama2026scud}, which carries a clone of the previous state in the
augmented vector. The lunar TDCP positioning and timekeeping results this
supports are those of \citet{iiyama2024navi}.

% ===========================================================================
\section{Visibility, Occultation and Channel Gating}
\label{sec:gnss-visibility}

Given a receiver state, \cpp{BuildChannels} walks every PRN of every attached
constellation and applies four gates in order. A PRN that fails any gate
contributes no rows to the measurement vector at that epoch, so the
measurement dimension changes from epoch to epoch.

\begin{figure}[htbp]
  \centering
  \begin{tikzpicture}[node distance=6mm and 9mm, >=Latex]
    \node[boxnode] (prn)  {for each PRN\\{\scriptsize of each constellation}};
    \node[boxnode, right=of prn] (fault) {fault gate\\{\scriptsize \code{IsFaultPrn}}};
    \node[boxnode, right=of fault] (build) {light time\\{\scriptsize \cpp{BuildChannelForPrn}}};
    \node[boxnode, below=of build] (vis) {occultation\\{\scriptsize \cpp{ComputeVisibility}}};
    \node[boxnode, left=of vis] (cn0) {$C/N_0$ gate\\{\scriptsize acquire / track}};
    \node[boxnode, left=of cn0] (sig) {noise sigmas\\{\scriptsize DLL / FLL / PLL}};

    \draw[flow] (prn) -- (fault);
    \draw[flow] (fault) -- (build);
    \draw[flow] (build) -- (vis);
    \draw[flow] (vis) -- (cn0);
    \draw[flow] (cn0) -- (sig);
    \draw[flow] (sig.south) -- ++(0,-0.55)
      node[boxnode, below, minimum width=3.1cm, anchor=north]
      {\code{std::vector<GnssChannel>}};
  \end{tikzpicture}
  \caption{Channel construction and gating inside
    \cpp{GNSSMeasurements::BuildChannels}. Each gate can drop the PRN; the
    surviving channels define the rows of the measurement vector at that
    epoch.}
  \label{fig:gnss-pipeline}
\end{figure}

\subsection{Fault gate}

PRNs registered with \cpp{SetFaultPrns} are skipped before any geometry is
computed. This is how a satellite outage or a clock-fault scenario is injected
for the fault-monitoring studies.

\subsection{Body-blocking test}
\label{sec:gnss-occultation}

Each entry of \cpp{SetOccludingBodies} is a sphere: a radius, and either a
fixed centre or a per-epoch position callback evaluated once per epoch (not
once per PRN). \cpp{ComputeVisibility}
(\file{cpp/lupnt/measurements/comms_utils.cc}) then runs two regimes.

\paragraph{Near-surface endpoints.} An endpoint with
$\lVert \vec{r} - \vec{r}_{\mathrm{body}} \rVert < R + h_{\min}$ (default
$h_{\min} = \SI{10}{\kilo\meter}$) is treated as a surface point, and the test
becomes an elevation mask. With $\vec{r}_{12} = \vec{r}_2 - \vec{r}_1$ and
$\vec{r}_{1b} = \vec{r}_{\mathrm{body}} - \vec{r}_1$,
\begin{equation}
  \label{eq:gnss-elevation}
  e = \cos^{-1}\!\left(
      \frac{\vec{r}_{12}\transpose \vec{r}_{1b}}
           {\lVert\vec{r}_{12}\rVert\,\lVert\vec{r}_{1b}\rVert}
    \right) - \frac{\pi}{2},
  \qquad
  \text{visible} \iff e > e_{\min} .
\end{equation}
$e$ is measured from the local horizon plane, positive upward. The default
mask is $e_{\min} = \SI{5}{\degree}$; the lunar-surface studies of
\citet[Ch.~6]{iiyama2026thesis} relax it to $\SI{-10}{\degree}$ to admit the
grazing rays that dominate a lunar user's sidelobe geometry.

\paragraph{Space-to-space links.} If neither endpoint is near the surface, the
test is a tangent-line occlusion test. With $\vec{r} = \vec{r}_2 - \vec{r}_1$
and $\vec{r}_{1b} = \vec{r}_1 - \vec{r}_{\mathrm{body}}$,
\begin{align}
  \label{eq:gnss-tangent}
  \theta_1 &= \cos^{-1}\!\left(
      \frac{-\vec{r}_{1b}\transpose\vec{r}}
           {\lVert\vec{r}\rVert\,\lVert\vec{r}_{1b}\rVert} \right),
  &
  \theta_2 &= \sin^{-1}\!\left(
      \frac{R}{\lVert\vec{r}_{1b}\rVert} \right),
  &
  d_{\mathrm{hor}} &= \sqrt{\lVert\vec{r}_{1b}\rVert^2 - R^2},
\end{align}
and the link is \emph{blocked} when
\begin{equation}
  \label{eq:gnss-occulted}
  \theta_1 < \theta_2
  \quad\text{and}\quad
  \lVert\vec{r}\rVert > d_{\mathrm{hor}} .
\end{equation}
$\theta_2$ is the angular radius of the body seen from $\vec{r}_1$ and
$\theta_1$ the angle between the line of sight and the direction to the body
centre, so the first condition says the ray points inside the body's angular
disc; the second says the far endpoint is beyond the tangent point rather than
in front of the body. Both are needed --- a receiver between the transmitter
and the body is not occulted even though the ray points at the body.

\begin{lstlisting}[language=cpp, caption={Space-to-space occultation test.},
label={lst:gnss-visibility}]
// cpp/lupnt/measurements/comms_utils.cc :: ComputeVisibility
Vec3 r = r2 - r1;
Real r_norm = r.norm();
Vec3 r1body = r1 - r_body;
Real dot_r_r1b = -r1body.dot(r);
Real theta1 = acos(clamp(dot_r_r1b / r_norm / r1body_norm, -1.0, 1.0));
Real theta2 = asin(clamp(R_body / r1body_norm, -1.0, 1.0));
Real r1_to_horizon = sqrt(r1body_norm * r1body_norm - R_body * R_body);
if (theta1 < theta2 && r_norm > r1_to_horizon) return false;

return true;
\end{lstlisting}

The body is a sphere with no atmosphere and no terrain: the test is purely
geometric, and an Earth-limb-grazing ray is declared visible with no
refraction, defocusing or atmospheric attenuation applied. For lunar users the
Moon is entered with its mean radius $R = \code{R\_MOON}$, so terrain relief of
a few kilometres is absorbed into the sphere radius.

\subsection{Boresight angles and the transmitter body frame}
\label{sec:gnss-attitude}

A surviving link still has to be strong enough to track, and the transmit gain
of a GNSS satellite varies by \SI{30}{\dB} between the main lobe and the
sidelobes a lunar receiver actually uses. The gain is a table lookup in the
satellite's yaw-steering body frame, so that frame must be built first.

\cpp{ComputeCN0} transforms the transmitter state, the receiver position and
the Sun position from \code{options.frame} into \code{Frame::GCRF} before
building the attitude, because ANTEX-derived patterns are defined in an
Earth-centred inertial frame. The orthonormal triad is
\begin{equation}
  \label{eq:gnss-attitude}
  \uvec{e}_z = \widehat{-\vec{r}_S}, \qquad
  \uvec{e}_y = \widehat{\uvec{e}_z \times \widehat{(\vec{r}_\odot - \vec{r}_S)}}, \qquad
  \uvec{e}_x = \widehat{\uvec{e}_y \times \uvec{e}_z},
\end{equation}
i.e.\ $\uvec{e}_z$ along nadir (the antenna boresight), $\uvec{e}_y$ along the
solar-panel axis and $\uvec{e}_x$ completing the right-handed triad. The
off-boresight and azimuth angles of the transmitter-to-receiver line of sight
$\uvec{u}$ are then
\begin{equation}
  \label{eq:gnss-boresight}
  \varphi_{\mathrm{tx}} = \cos^{-1}\!\bigl(\uvec{u}\transpose\uvec{e}_z\bigr),
  \qquad
  \theta_{\mathrm{tx}} = \atantwo\!\bigl(\uvec{u}\transpose\uvec{e}_y,\ \uvec{u}\transpose\uvec{e}_x\bigr),
\end{equation}
and the receiver off-boresight angle, measured from the receiver's pointing
target $\vec{b}$ (the Moon or Earth centre, supplied by
\cpp{SetBoresightTargetProvider}), is
\begin{equation}
  \label{eq:gnss-boresight-rx}
  \varphi_{\mathrm{rx}}
  = \cos^{-1}\!\left(
      \widehat{(\vec{b} - \vec{r}_R)}\transpose\,\widehat{(\vec{r}_S - \vec{r}_R)}
    \right).
\end{equation}

The velocity-aware overload derives the same frame from the documented nominal
yaw-steering law,
\begin{equation}
  \label{eq:gnss-yaw}
  \psi_{\mathrm{nom}} = \atantwo\bigl(-\tan\beta,\ \sin\mu\bigr),
\end{equation}
with $\beta$ the Sun elevation above the orbital plane and $\mu$ the orbit
angle from the midnight point, then rotates the orbital reference frame
($\uvec{e}_x^{\mathrm{ref}} = \widehat{\vec{v}_S}$) about nadir by
$\psi_{\mathrm{nom}}$. The two constructions are numerically identical, since
the nominal yaw law is by definition the yaw that keeps the panels
Sun-pointing \citep{kouba2009yaw,cheng2025yaw}.

\begin{lstlisting}[language=cpp, caption={Body frame, boresight angles and
link budget for one channel.}, label={lst:gnss-cn0}]
// cpp/lupnt/measurements/gnss_measurement.cc :: GNSSMeasurements::ComputeCN0
GnssAttitude::Compute(rv_tx_gcrf.head(3), Vec3(rv_tx_gcrf.tail(3)), r_sun_gcrf, ex, ey, ez);
Real phi_tx = safe_acos(u_tx2rx_gcrf.dot(ez));
Real theta_tx = atan2(u_tx2rx_gcrf.dot(ey), u_tx2rx_gcrf.dot(ex));
Real phi_rx = safe_acos(u_rx2body.dot(u_rx2tx));

Real G_tx = constellation->GetTransmitterAntenna(channel.prn, channel.frequency)
                .ComputeGain(theta_tx, phi_tx);
Real G_rx = rx_antenna_.ComputeGain(0.0, phi_rx);
Real P_tx = constellation->GetTransmitPowerDbw(channel.prn, channel.frequency);

return LinkBudget(P_tx, G_tx, G_rx, range, channel.frequency, rx_params_);
\end{lstlisting}

Setting \code{options.tx\_yaw\_model = TxYawModel::DEDICATED} replaces
\cref{eq:gnss-yaw} with the block-specific eclipse yaw law of
\cpp{GnssYawSteering}
(\cpp{Gps3EclipseYawAngle} for GPS,
\cpp{GalileoIovEclipseYawAngle} for Galileo, nominal otherwise), evaluated
continuously from geometry. These laws differ from nominal only near the
orbit noon/midnight singularity at small $\beta$, where the real satellite
rate-limits its yaw; because there is no maneuver-window state machine, the
result approximates the eclipse-season attitude rather than reproducing a
maneuver schedule. \citet[Ch.~6.4.3]{iiyama2026thesis} assumes nominal yaw
steering throughout.

\subsection{Antenna gain gating}
\label{sec:gnss-gain}

\cpp{Antenna::ComputeGain} (\file{cpp/lupnt/measurements/antenna.cc})
interpolates the measured pattern, wrapping $\theta$ to $[0,360]^\circ$ and
$\varphi$ to $[-180,180]^\circ$:

\begin{lstlisting}[language=cpp, caption={Pattern lookup, with the implicit
sidelobe cutoff.}, label={lst:gnss-antenna}]
// cpp/lupnt/measurements/antenna.cc :: Antenna::ComputeGain
Real gain = 0.0;
if (n_dim_ == 0) return gain;                       // omni antenna: G == 0 dB
theta = WrapToTwoPi(theta) * DEG;
phi = WrapToPi(phi) * DEG;
if (phi > phi_max_ || phi < -phi_max_) return NAN;  // off-pattern -> NaN
if (phi < 0 && phi_(0) >= 0) phi = -phi;            // symmetry
if (n_dim_ == 2) gain = LinearInterp2d(phi_, theta_, gain_, phi.val(), theta.val());
else if (n_dim_ == 1) gain = LinearInterp1d(phi_, gain_, phi.val());
return gain;
\end{lstlisting}

An empty antenna name gives an omnidirectional pattern ($G \equiv 0$
\si{\dB}), and a direction outside the tabulated $\varphi$ range returns
\code{NaN}. The \code{NaN} propagates through the link budget, and the
gate below drops the channel, so the pattern's angular extent acts as an
implicit sidelobe cutoff. The same happens when
\cpp{HasTransmitterInfo} is false for a PRN/frequency pair: $C/N_0$ is
\code{NaN} and the channel is dropped.

The received carrier-to-noise density follows the closed-form budget
\citep[Eq.~6.59]{iiyama2026thesis}
\begin{equation}
  \label{eq:gnss-cn0}
  \left(\frac{C}{N_0}\right)_{\si{\dBHz}}
  = P_{\mathrm{tx}}
  + G_{\mathrm{tx}}(\theta_{\mathrm{tx}}, \varphi_{\mathrm{tx}})
  + G_{\mathrm{rx}}(\varphi_{\mathrm{rx}})
  - L_{\mathrm{fs}} - L_{\mathrm{atm}} - L_{\mathrm{ad}} - L_{\mathrm{pol}}
  - 10\log_{10} k_B - 10\log_{10} T_{\mathrm{eff}} ,
\end{equation}
with $L_{\mathrm{fs}} = 20\log_{10}(4\pi \rho f / c)$ the free-space path
loss. The EIRP values, the measured gain patterns and the mapping from
$C/N_0$ to the DLL, FLL and PLL noise standard deviations are specified in
\cref{chap:link-budget}; here they matter only through the gate and through
the diagonal of the measurement covariance.

\subsection{Acquisition and tracking hysteresis}

The $C/N_0$ gate is stateful. A PRN not currently in lock must exceed the
acquisition threshold; one already in lock need only exceed the (lower)
tracking threshold:
\begin{equation}
  \label{eq:gnss-cn0-gate}
  \text{keep} \iff
  \frac{C}{N_0} >
  \begin{cases}
    (C/N_0)_{\mathrm{acq}} = \SI{22}{\dBHz}, & \text{PRN not in lock},\\[2pt]
    (C/N_0)_{\mathrm{trk}} = \SI{20}{\dBHz}, & \text{PRN in lock}.
  \end{cases}
\end{equation}
The lock set \code{tracking\_prns\_} is rebuilt after every
\cpp{BuildChannels} call from the surviving channels, and is cleared by
\cpp{ResetTracking}. The \SI{2}{\dB} hysteresis reproduces the observed
behaviour of a real receiver that holds a weak satellite it has already
acquired, and prevents the rapid acquire/drop chatter that a single threshold
produces at the sidelobe boundaries. The comparison is written
\code{!(cn0 > threshold)}, so a \code{NaN} $C/N_0$ is dropped rather than
kept. Setting \code{options.apply\_cn0\_threshold = false} disables the gate
entirely, and \cpp{SetCN0Threshold} collapses both thresholds to one value
(the deprecated single-threshold behaviour).

% ===========================================================================
\section{Measurement Jacobians}
\label{sec:gnss-jacobian}

\cpp{GnssMeasurement::ComputeVector} returns $\vec{y} = h(\vec{x})$ and, when
asked, the analytic Jacobian $\mat{H} = \partial h/\partial\vec{x}$ with
$|\code{observables}|$ rows and $\dim\vec{x}$ columns, zero-initialised so
that any state component not named in \code{indices} keeps a zero column.

\subsection{Building blocks}

Differentiating \cref{eq:gnss-range,eq:gnss-rangerate} with respect to the
receiver position and velocity, holding the transmitter state fixed,
\begin{align}
  \label{eq:gnss-drho}
  \frac{\partial \rho}{\partial \vec{r}_R} &= \uvec{u}\transpose,
  &
  \frac{\partial \rho}{\partial \vec{v}_R} &= \vec{0}\transpose, \\
  \label{eq:gnss-drhodot}
  \frac{\partial \dot\rho}{\partial \vec{r}_R}
    &= \left( \frac{\Delta\vec{v}}{\rho}
       - \frac{\dot\rho\,\Delta\vec{r}}{\rho^{2}} \right)^{\!\mathsf{T}},
  &
  \frac{\partial \dot\rho}{\partial \vec{v}_R} &= \uvec{u}\transpose .
\end{align}
\cref{eq:gnss-drhodot} can equivalently be written
\begin{equation}
  \label{eq:gnss-drhodot-proj}
  \frac{\partial \dot\rho}{\partial \vec{r}_R}
  = \frac{1}{\rho}
    \left[ \bigl( \mat{I} - \uvec{u}\,\uvec{u}\transpose \bigr) \Delta\vec{v}
    \right]^{\mathsf{T}},
\end{equation}
which makes the structure explicit: only the component of the relative
velocity \emph{perpendicular} to the line of sight changes the range rate when
the receiver moves, and the sensitivity falls off as $1/\rho$.

The clock scale factors are the derivatives of \cref{eq:gnss-clock-units},
\begin{equation}
  \label{eq:gnss-clock-deriv}
  \frac{\partial \Delta\rho_R}{\partial b_R}
    = c\,\frac{\partial \Delta t_R}{\partial b_R}
    = \frac{c}{s(\text{unit})},
  \qquad
  \frac{\partial \Delta\dot\rho_R}{\partial d_R}
    = \frac{c}{s(\text{unit})},
\end{equation}
where $s(\text{unit})$ is the seconds-to-unit scale: $s = 1$ for
\code{SECONDS}, so the factor is $c$; $s = 1/c$ for \code{METERS}, so the
factor is $1$; and $s = \num{e-3}/c$ for \code{KILOMETERS}.

\subsection{Rows}

\begin{equation}
  \label{eq:gnss-jac-pseudorange}
  \frac{\partial P}{\partial \vec{r}_R} = \uvec{u}\transpose,
  \qquad
  \frac{\partial P}{\partial b_R} = \frac{c}{s},
  \qquad
  \frac{\partial P}{\partial \vec{v}_R} = \vec{0}\transpose,
  \qquad
  \frac{\partial P}{\partial d_R} = 0 ;
\end{equation}
\begin{equation}
  \label{eq:gnss-jac-doppler}
  \frac{\partial D}{\partial \vec{r}_R}
    = -\frac{1}{\lambda}\left( \frac{\Delta\vec{v}}{\rho}
       - \frac{\dot\rho\,\Delta\vec{r}}{\rho^{2}} \right)^{\!\mathsf{T}},
  \qquad
  \frac{\partial D}{\partial \vec{v}_R} = -\frac{1}{\lambda}\uvec{u}\transpose,
  \qquad
  \frac{\partial D}{\partial d_R} = -\frac{c}{\lambda s} ;
\end{equation}
\begin{equation}
  \label{eq:gnss-jac-carrier}
  \frac{\partial \Phi}{\partial \vec{r}_R} = \frac{1}{\lambda}\uvec{u}\transpose,
  \qquad
  \frac{\partial \Phi}{\partial b_R} = \frac{c}{\lambda s},
  \qquad
  \frac{\partial \Phi_{\mathrm{meas}}}{\partial N} = 1 .
\end{equation}

For the default eight-state layout $\vec{x} = [\vec{r}_R;\ \vec{v}_R;\ b_R;\
d_R]$ with the clock in seconds, the full single-channel Jacobian is
\begin{equation}
  \label{eq:gnss-jac-block}
  \mat{H} =
  \begin{bmatrix}
    \uvec{u}\transpose & \vec{0}\transpose & c & 0 \\[4pt]
    -\frac{1}{\lambda}\left(\frac{\Delta\vec{v}}{\rho}
      - \frac{\dot\rho\,\Delta\vec{r}}{\rho^{2}}\right)^{\!\mathsf{T}}
      & -\frac{1}{\lambda}\uvec{u}\transpose & 0 & -\frac{c}{\lambda} \\[6pt]
    \frac{1}{\lambda}\uvec{u}\transpose & \vec{0}\transpose & \frac{c}{\lambda} & 0
  \end{bmatrix},
\end{equation}
and \cpp{GNSSMeasurements::ComputeFromChannels} stacks one such block per
surviving channel, giving a $3 n_{\mathrm{ch}} \times \dim\vec{x}$ matrix at
that epoch.

\begin{lstlisting}[language=cpp, caption={Analytic Jacobian rows.},
label={lst:gnss-jacobian}]
// cpp/lupnt/measurements/gnss_measurement.cc :: GnssMeasurement::ComputeVector
Vec3 u = dr / rho;
Vec3 d_rhodot_dr = dv / rho - rho_dot * dr / (rho * rho);
Vec3 d_rhodot_dv = u;
Real db_m  = ClockBiasMetersDerivative(options.clock_bias_unit);
Real dd_mps = ClockDriftMetersPerSecondDerivative(options.clock_bias_unit);

case GnssObservable::PSEUDORANGE:
  H->block(row, idx.position, 1, 3) = ToRowVec3d(u);
  if (idx.clock_bias >= 0 && idx.clock_bias < user_state.size())
    (*H)(row, idx.clock_bias) = db_m.val();
  break;

case GnssObservable::DOPPLER:
  H->block(row, idx.position, 1, 3) = ToRowVec3d(-d_rhodot_dr / lambda);
  H->block(row, idx.velocity, 1, 3) = ToRowVec3d(-d_rhodot_dv / lambda);
  if (idx.clock_drift >= 0 && idx.clock_drift < user_state.size())
    (*H)(row, idx.clock_drift) = (-dd_mps / lambda).val();
  break;

case GnssObservable::CARRIER_PHASE:
  H->block(row, idx.position, 1, 3) = ToRowVec3d(u / lambda);
  if (idx.clock_bias >= 0 && idx.clock_bias < user_state.size())
    (*H)(row, idx.clock_bias) = (db_m / lambda).val();
  if (idx.carrier_integer >= 0 && idx.carrier_integer < user_state.size())
    (*H)(row, idx.carrier_integer) = 1.0;
  break;
\end{lstlisting}

\subsection{What the Jacobian omits}

\cref{eq:gnss-jac-block} is a \emph{partial} derivative that holds the
transmitter state, the delays and the noise model fixed. Omitted are:

\begin{itemize}[leftmargin=1.4em,itemsep=1pt]
  \item the light-time coupling
    $\partial \vec{r}_S / \partial \vec{r}_R
     = \vec{v}_S\,\uvec{u}\transpose / c$, of relative size
    $\lVert\vec{v}_S\rVert/c \approx \num{1.3e-5}$;
  \item $\partial \Delta\rho_{\mathrm{Sh}} / \partial \vec{r}_R$, of order
    \num{1e-9};
  \item $\partial \Delta\rho_{\mathrm{pl}} / \partial \vec{r}_R$, which is not
    even available in closed form for the ray-traced model;
  \item the dependence of $C/N_0$, and hence of the measurement covariance
    $\mat{R}$, on the receiver state --- a genuine effect near a gain-pattern
    null, but one that a Kalman filter is not structured to use anyway;
  \item any derivative with respect to the transmitter clock, since the
    transmitter clock is not estimated.
\end{itemize}

All of these are far below the linearisation error of an extended Kalman
filter at the position uncertainties of interest. When exact derivatives are
required, the whole model is written in the forward-mode \code{Real} scalar
type and can be differentiated end to end by automatic differentiation
\citep{leal2018autodiff} instead of using \cref{eq:gnss-jac-block}.

\subsection{Measurement covariance}

The diagonal covariance is built from the per-channel sigmas
$\sigma_P$, $\sigma_D$, $\sigma_\Phi$ that
\cpp{BuildChannels} derived from $C/N_0$ through the DLL, FLL and PLL noise
formulas of \cref{chap:link-budget}:
\begin{equation}
  \label{eq:gnss-covariance}
  \mat{R} = \diag\bigl(
    \sigma_{P,1}^2,\ \sigma_{D,1}^2,\ \sigma_{\Phi,1}^2,\ \ldots,\
    \sigma_{P,n}^2,\ \sigma_{D,n}^2,\ \sigma_{\Phi,n}^2 \bigr).
\end{equation}
A sigma that is not finite and positive is replaced by $1$ in the observable's
own unit, so a channel with no link-budget information contributes a
unit-weight row rather than a singular covariance. The Doppler and
carrier-phase sigmas are stored in \si{\hertz} and cycles respectively,
obtained by dividing the range-rate and phase sigmas by $\lambda$. Correlation
between observables of the same channel, and between channels, is not
modelled.

% ===========================================================================
\section{Online and Precomputed Evaluation}
\label{sec:gnss-precompute}

Two evaluation paths must agree. The online path computes
$\mathcal{Y}(t_i, \vec{x}_i) = \cpp{Compute}(t_i, \vec{x}_i)$ one epoch at a
time, and is what a filter calls through \cpp{CreateFunction}. The precompute
path evaluates the same epochs for vectors \code{receive\_times} and
\code{user\_states} at once.

They differ only when a \emph{batch} plasma provider is installed. A ray-traced
delay is far too expensive to evaluate inside a filter update, and a batch
provider may want the whole epoch/channel grid at once (for example to
vectorise the ray tracing, or to farm it out to a separate process). In that
case \cpp{Precompute} runs in three ordered stages:

\begin{enumerate}[leftmargin=1.6em,itemsep=1pt]
  \item build all channels for all epochs with
    \code{compute\_ionosphere\_plasma\_delay = false}, so visibility, light
    time and the transmitter clock terms are resolved but
    $\Delta\rho_{\mathrm{pl}} = 0$;
  \item call the batch provider once with
    (\code{receive\_times}, \code{user\_states}, \code{channels\_by\_epoch})
    and write the returned $\Delta\rho_{\mathrm{pl},ij}$ into each channel,
    with the epoch and channel counts checked;
  \item evaluate the observables from the patched channels.
\end{enumerate}

This ordering is what avoids an online ray trace whose result would only be
overwritten by the batch result.

\begin{lstlisting}[language=cpp, caption={Batch plasma staging inside
\cpp{Precompute}.}, label={lst:gnss-precompute}]
// cpp/lupnt/measurements/gnss_measurement.cc :: GNSSMeasurements::Precompute
if (options_.apply_ionosphere_plasma_delay && batch_custom_ionosphere_plasma_delay_model_) {
  std::vector<std::vector<GnssChannel>> channels_by_epoch;
  for (size_t i = 0; i < receive_times.size(); i++)
    channels_by_epoch.push_back(BuildChannels(receive_times[i], user_states[i], false));

  std::vector<std::vector<Real>> delays_m = batch_custom_ionosphere_plasma_delay_model_(
      receive_times, user_states, channels_by_epoch);

  for (size_t i = 0; i < channels_by_epoch.size(); i++) {
    for (size_t j = 0; j < channels_by_epoch[i].size(); j++)
      channels_by_epoch[i][j].ionosphere_plasma_delay_m = delays_m[i][j];
    MatXd H;
    precomputed_.push_back(ComputeFromChannels(receive_times[i], user_states[i],
                                               channels_by_epoch[i],
                                               compute_jacobians ? &H : nullptr));
  }
  return precomputed_;
}
\end{lstlisting}

In the staged Lunar GNSS ODTS scenario the same contract is realised across a
file boundary: \cpp{LunarGnssOdtsApp::Precompute} writes all
light-time-corrected links with zero plasma delay, the Python driver
\file{python/examples/ex6_precompute.py} fills the delay terms with GCPM/IRI
ray tracing (\cref{chap:ionosphere}), and \cpp{LunarGnssOdtsApp::Step} merges
the delay
file back into the truth channels before generating pseudorange, Doppler and
optional TDCP measurements.

\cref{tab:gnss-options} lists the configuration knobs that change the model's
behaviour.

\begin{table}[htbp]
  \centering
  \caption{\cpp{GnssMeasurementOptions} --- the knobs that change the model.}
  \label{tab:gnss-options}
  \footnotesize
  \begin{tabularx}{\textwidth}{p{0.36\textwidth}lL}
    \toprule
    Option & Default & Effect \\
    \midrule
    \code{observables} & all three & selects and orders the emitted rows \\
    \code{indices} & $0,3,6,7,-1$ & state layout; a negative index drops the term \\
    \code{clock\_bias\_unit} & \code{SECONDS} & unit of $b_R$, $d_R$; sets the Jacobian scale \\
    \code{receive\_time\_scale} & \code{TAI} & scale of the epochs passed in \\
    \code{ephemeris\_time\_scale} & \code{TAI} & must be \code{TAI}; asserted \\
    \code{frame} & \code{ECI} & common frame of receiver and transmitter \\
    \code{solve\_light\_time} & true & enables \cref{eq:gnss-lighttime-offset} \\
    \code{recompute\_transmit\_time}\newline\code{\_from\_ephemeris} & false & re-solve light time per state evaluation \\
    \code{light\_time\_max\_iterations} & 10 & iteration cap \\
    \code{light\_time\_tolerance\_s} & \num{1e-11} & $\epsilon_{\mathrm{lt}}$ of \cref{eq:gnss-lighttime-stop} \\
    \code{apply\_transmitter\_relativity} & true & enables \cref{eq:gnss-relativity} \\
    \code{apply\_shapiro\_delay} & true & enables \cref{eq:gnss-shapiro} \\
    \code{apply\_visibility} & true & enables the occultation gate \\
    \code{apply\_cn0\_threshold} & true & enables \cref{eq:gnss-cn0-gate} \\
    \code{cn0\_acquisition\_threshold\_dbhz} & \num{22} & acquire threshold [\si{\dBHz}] \\
    \code{cn0\_tracking\_threshold\_dbhz} & \num{20} & hold threshold [\si{\dBHz}] \\
    \code{apply\_ionosphere\_plasma\_delay} & false & enables \cref{eq:gnss-plasma} \\
    \code{default\_ionosphere\_plasma\_delay\_m} & \num{0} & constant fallback delay \\
    \code{tx\_yaw\_model} & \code{NOMINAL} & \code{NOMINAL} or \code{DEDICATED} eclipse yaw \\
    \bottomrule
  \end{tabularx}
\end{table}

% ===========================================================================
\section{Model Boundaries}
\label{sec:gnss-boundaries}

\subsection{What is approximated}

\begin{description}[leftmargin=1.4em,style=nextline,font=\normalfont\itshape]
  \item[Transmitter ephemeris interpolation]
    In \cpp{FitStateHistoryChebyshev} the constellation-level fit is currently
    degree one over one tabulation step, i.e.\ piecewise-linear in practice.
    The
    chord error is $\tfrac{1}{8}\lVert\ddot{\vec{r}}_S\rVert (\Delta t)^2$;
    for a GPS satellite ($\lVert\ddot{\vec{r}}_S\rVert \approx
    \SI{0.56}{\meter\per\second\squared}$) this is about \SI{0.7}{\meter} at
    $\Delta t = \SI{10}{\second}$ and \SI{70}{\meter} at
    $\Delta t = \SI{100}{\second}$. The tabulation grid, not the fit, sets the
    accuracy. The SP3 loader's own fit is different and much better: degree up
    to 11 over eight steps, sampled by a tenth-order Lagrange interpolant,
    with velocity as the exact analytic derivative.
  \item[Light time]
    One-way, geometric, iterated to \SI{1e-11}{\second} $\approx
    \SI{3}{\milli\meter}$. Solved in the ephemeris frame with no
    frame-rotation term, which is exact for an inertial frame and would not be
    for an Earth-fixed one.
  \item[Relativity]
    The transmitter contribution is the periodic
    eccentricity term \cref{eq:gnss-relativity} plus the Sun's Shapiro delay.
    The constant fractional-rate offset of the satellite oscillator is assumed
    to be already absorbed into the broadcast/precise clock, per
    \citet{icdgps200}. Earth's Shapiro delay, second-order Doppler on the
    receiver clock, and the frame-dependence noted in
    \cref{sec:gnss-clocks} are not modelled here; see
    \cref{chap:time-conversions} for the receiver-side timescale treatment.
  \item[Occulting bodies]
    Spheres with a single radius and an optional elevation mask. No terrain,
    no atmosphere, no refraction, no limb attenuation, no penumbra.
  \item[Transmitter attitude]
    Nominal yaw steering, with an optional continuously-evaluated
    block-specific eclipse law that is not a maneuver schedule. Antenna phase
    variations (PCV) are not applied --- only the constant phase-centre offset
    (\cref{eq:gnss-pco}) is.
  \item[Noise]
    Diagonal, derived from $C/N_0$ through steady-state tracking-loop
    formulas. No multipath, no inter-channel or inter-observable correlation,
    no time correlation, no cycle slips.
\end{description}

\subsection{What is omitted}

\begin{itemize}[leftmargin=1.4em,itemsep=1pt]
  \item \textbf{Carrier phase wind-up} (\cref{sec:gnss-windup}). Absorbed into
    $\phi_0$ and the ambiguity; up to one cycle per receiver revolution.
  \item \textbf{Troposphere} in the GNSS path (\cref{sec:gnss-tropo}). The
    Saastamoinen model exists but is wired only into the ground-station
    application.
  \item \textbf{Group delay}. The \code{group\_delay\_s} field is honoured by
    \cref{eq:gnss-tx-clock} but is never populated from the navigation
    message; it defaults to zero. Single-frequency users of a real broadcast
    clock therefore inherit a systematic bias of a few nanoseconds.
  \item \textbf{Solid-Earth and ocean loading, antenna PCV, receiver
    hardware delays.} Not applicable to a space receiver, or absorbed into
    $\phi_0$.
  \item \textbf{Time derivatives of the medium delays} in the Doppler
    observable.
  \item \textbf{Multipath and near-field effects.} No surface model is coupled
    to the measurement path.
  \item \textbf{GLONASS.} Its FDMA frequency plan and tabulated-state
    navigation message are not supported by \cpp{RinexNavLoader}.
  \item \textbf{Galileo GST/GPST (GAGP) correction}, worth well under
    \SI{100}{\nano\second} of satellite clock.
\end{itemize}

\subsection{Expected accuracy}

With a precise (SP3 plus ANTEX) constellation, a dense ephemeris grid, Shapiro
enabled and plasma supplied by the ray tracer of \cref{chap:ionosphere}, the
modelled pseudorange reproduces the geometric truth at the level of the
ephemeris interpolation error --- decimetres for a
\SIrange{10}{30}{\second} grid --- plus the plasma model error, which
dominates at low geocentric altitude and is characterised in
\cref{chap:ionosphere}. Carrier phase is consistent with pseudorange to the
same level, with the code--carrier divergence
$2\Delta\rho_{\mathrm{pl}}$ correctly signed. Doppler carries an additional
unmodelled $\dd\Delta\rho_{\mathrm{pl}}/\dd t$ of order
\SI{0.01}{\meter\per\second}.

With a broadcast (BRDC) constellation the error is deliberately larger and is
the point of the exercise: it is the real signal-in-space error, roughly
\SIrange{0.5}{2}{\meter} in orbit and \SIrange{1}{3}{\meter} in clock, and
that is what the filter must be robust to. The synthetic SISE of
\cref{eq:gnss-sise} substitutes for it at future epochs with per-PRN
constant draws rather than the true time-correlated process, which is
optimistic over arcs longer than a few hours.

The navigation performance these numbers support --- lunar orbit
determination at the tens-of-metres level and clock estimation at the
nanosecond level from terrestrial GNSS sidelobes --- is reported in
\citet{iiyama2024navi} and \citet{iiyama2026scud}, and analysed in full in
\citet[Ch.~6]{iiyama2026thesis}.
